\section{Implementing Earley's algorithm in F\#}
The chart algorithm above is implemented directly using a record for an Earley item.
\texttt{Rule} identifies its production, \texttt{Dot} its recognized prefix length and \texttt{Origin} its start position.
Immutable sets prevent duplicate items. The chart is an array whose entries are replaced as prediction and completion
enlarge each finite set. These representations correspond to the lecture's $Q_0,\ldots,Q_n$ without hiding
state inside a parser object.
@@snippet:earley-items@@
@@snippet:earley-closure@@

At every position, prediction and completion repeat until neither changes the current chart set.
Only then does scanning move matching terminal items into the next set. This ordering also handles nullable productions:
an item can be predicted and completed without advancing the input position.
\href{../../Chapter-08/Earley-Parser-Pure.ipynb}{\nolinkurl{Chapter-08/Earley-Parser-Pure.ipynb}}
checks recognition; \href{../../Chapter-08/Earley-Parser.ipynb}{\nolinkurl{Chapter-08/Earley-Parser.ipynb}}
constructs parse trees from completed spans. Tree construction records an active span path to avoid following a
cyclic nullable derivation forever. The recognition algorithm and tree enumeration remain separate responsibilities.
